\section{AI、能源与文明尺度}

前一章把商业化落到电站成本和燃料路线，本章把视角拉回这期访谈开头的 AI 与文明问题。为什么一个核聚变创业者会谈 AGI、数据中心和算力？因为能源不是下游应用，而是所有计算、制造和社会组织的底座；当 AI 继续指数增长时，能源供给会从背景变量变成前台约束。

本章把核聚变放回 AI 时代。杨钊认为，任何指数增长的东西都会遇到瓶颈；AI 的短期瓶颈可能是算力、芯片、数据和能源，长期看能源供给会成为真正的大瓶颈。只要能提供更便宜的能源，需求几乎一定会被用掉，就像计算性能一旦提升，总会有新的应用把它消耗掉。

反过来，AI 也能帮助核聚变。访谈中提到三类作用：第一，AI 可以用于实时控制，替代复杂但太慢的传统模型；第二，AI 可以增强诊断，用算法替代部分昂贵硬件；第三，AI 可以用真实实验数据训练装置级仿真模型，缩短实验迭代周期。也就是说，AI 和聚变不是单向关系：AI 需要能源，聚变也需要 AI 降本增效。

\begin{figure}[H]
\centering
\includegraphics[width=0.94\textwidth]{energy-ai-loop.png}
\caption{AI 和能源互相放大：算力增长推高电力需求，低成本能源反过来抬高文明上限。自制概念图，依据 01:44:42--01:50:00 对谈内容整理。}
\end{figure}

\begin{knowledgebox}{读图：能源不是 AI 的外部背景}
图中不是简单地把 Fusion 接到 Data Center，而是显示一个正反馈：AI 消耗更多能源，推动新能源投资；新能源降低成本，又让更多计算、制造和实验成为可能。核聚变如果成功，会成为 AI 时代的底层放大器之一。
\end{knowledgebox}

\begin{knowledgebox}{AI 对聚变的三类帮助}
\begin{tabularx}{\textwidth}{p{0.20\textwidth}p{0.34\textwidth}X}
\toprule
场景 & 做什么 & 价值 \\
\midrule
实时控制 & 用 AI 近似复杂物理过程，快速给出控制策略 & 缩短控制延迟，提高等离子体稳定性。 \\
诊断增强 & 用算法提升观测精度或替代昂贵诊断设备 & 降低硬件成本，提升状态估计质量。 \\
实验仿真 & 用真实实验数据训练装置级 surrogate model & 减少试错次数，加快调参和运行优化。 \\
\bottomrule
\end{tabularx}
\end{knowledgebox}

\subsection{安全、监管与燃料选择}

本节补上商业电站必须面对的监管问题。上一节说 AI 可以降本增效，但能源产品不是实验室模型，燃料、废料和监管会直接变成成本。杨钊特别区分氘氚和氘氘，是因为它们的物理难度与商业难度方向相反：氘氚更容易反应，氘氘更接近长期低监管、低燃料压力的愿景。

聚变电站不仅要 Q 高，还要可监管、可维护、可被社会接受。访谈中特别提到氘氚与氘氘路线差异。氚可用于武器相关体系，监管强、成本高，而且自然界难以稳定大量存在，需要边发电边增殖氚；这会增加工程复杂度和监管成本。能量奇点长期希望做氘氘聚变，因为氘在海水中储量极大，原料可得性、安全性和废料处理更有优势，但它需要更高参数。

\begin{warningbox}{燃料路线影响的不只是物理难度}
氘氚更容易达到反应条件，但氚的增殖、监管和安全系统会增加电站复杂度；氘氘更难，但若能实现，原料和监管压力可能显著下降。商业化路线要同时比较物理难度和整站成本。
\end{warningbox}

\begin{knowledgebox}{燃料路线对照}
\begin{tabularx}{\textwidth}{p{0.18\textwidth}p{0.34\textwidth}X}
\toprule
路线 & 优势 & 难点 \\
\midrule
氘氚 & 更容易达到反应条件，是近期主流验证路线 & 氚受监管、需增殖、系统安全和成本压力高。 \\
氘氘 & 氘来自海水，长期原料和监管优势明显 & 反应更难，需要更高三重积和更强装置性能。 \\
示范电站 & 可先证明 Q、取热和并网可行 & 还不等于最低成本、最低监管的终局路线。 \\
\bottomrule
\end{tabularx}
\end{knowledgebox}

\subsection{无限能源之后会怎样}

本节把讨论从电站经济性推到文明尺度。前面的每一步都很工程化：磁体、Q、成本、监管、融资；但杨钊之所以愿意把这件事当作“有生之年系列”，是因为它一旦跨过商业化阈值，就不只是在替代某一种电源，而是在扩大人类可以承担的能源预算。

当能源成本大幅下降，很多今天“不划算”的事情会变得可行：海水淡化、大规模计算、材料制造、极端环境工程、行星际推进等都会获得新边界。杨钊把聚变看作文明尺度的变化：它不只是替代火电，而是可能让人类使用能源的方式整体跨一个数量级。

\begin{importantbox}{能源供给的文明含义}
能源不是一个单独行业的输入，而是计算、制造、交通、城市、农业和太空活动的底层约束。聚变如果把低成本稳定供能做成，它改变的是文明可做事情的集合。
\end{importantbox}

\begin{warningbox}{不要把“无限能源”理解成没有约束}
即使聚变成功，设备建造、材料、运维、监管、输电、电网调度和资本周期仍然存在。更准确的说法是：聚变可能把能源边际成本和资源约束显著下降，但不会让所有工程约束自动消失。
\end{warningbox}

\subsection{失败风险与爬山模型}

上一节讲的是能源成功之后的文明空间，本节回到创业者每天面对的现实问题：这座山怎么爬，在哪里可能摔下来。核聚变的宏大愿景很容易让人忽略执行层面的残酷性，但杨钊的结尾判断很清楚：只要有足够时间和钱，路线在工程上有解；真正的风险是启动不了、交付不了，或者团队无法持续升级自己的能力。

访谈尾声里，杨钊把聚变创业比作持续爬山，而且山的高度是指数增长的。失败原因可以被压缩成三类：钱不够，下一台装置启动不了；人不够，关键工程判断和执行跟不上；事没做成，装置没有达到设计目标，后续 380 就卖不出去。这个判断很冷静：聚变的 PMF 并不模糊，需求是真需求，产品也明确，真正风险在于能不能把产品做出来。

\begin{knowledgebox}{风险拆解：为什么不是普通创业}
\begin{tabularx}{\textwidth}{p{0.18\textwidth}p{0.34\textwidth}X}
\toprule
风险 & 具体表现 & 后果 \\
\midrule
钱不够 & 170 或关键子系统无法启动 & 技术路线只能停在前一座山。 \\
人不够 & 无法在信息不足时做正确工程判断 & 进度拖慢，错误难定位。 \\
事没做成 & 交付参数达不到设计目标 & 380 和商业电站销售逻辑失效。 \\
\bottomrule
\end{tabularx}
\end{knowledgebox}

\subsection{本章小结}

AI 的长期增长需要能源，核聚变的研发也可被 AI 加速。聚变商业化的意义不仅是新增发电方式，而是可能改变计算、工业和文明扩张的边界。

